\documentclass[11pt,a4paper]{article}
\usepackage[margin=1in]{geometry}
\usepackage[T1]{fontenc}
\usepackage[utf8]{inputenc}
\usepackage{lmodern}
\usepackage{hyperref}
\usepackage{longtable}
\usepackage{array}
\usepackage{enumitem}
\setlist[itemize]{leftmargin=1.2em}
\setlength{\parskip}{0.5em}
\setlength{\parindent}{0pt}

\title{Domain Body Test Case: In-Depth Codebase Study Guide}
\author{Generated for learning and MATLAB replication}
\date{\today}

\begin{document}
\maketitle

\section*{1. Big Picture}
This repository implements a distributed automotive feature simulation using SIL Kit.
It models five logical ECU participants:
\begin{itemize}
  \item Plant Model RestBus
  \item Zonal A
  \item HPC
  \item Zonal B
  \item Telematics ECU
\end{itemize}

Core feature under test: trunk open/close behavior under ignition and vehicle speed conditions, with hazard forwarding in parallel.

Signal flow summary:
\begin{itemize}
  \item PlantModel sends CAN stimulus to Zonal A (CAN1).
  \item Zonal A wraps CAN payload into Ethernet and sends to HPC (Eth4).
  \item HPC applies logic, then sends:
  \begin{itemize}
    \item hazard-related Ethernet to Zonal B (Eth2),
    \item trunk-related Ethernet to Zonal B (Eth3),
    \item trunk response Ethernet to Telematics (Eth1).
  \end{itemize}
  \item Zonal B converts incoming Ethernet payloads back to CAN (CAN2/CAN3) to PlantModel.
  \item PlantModel validates expected payloads and prints PASS/FAIL.
\end{itemize}

\section*{2. Repository Structure}
\textbf{Application logic folders}
\begin{itemize}
  \item \texttt{Zonal\_A/}
  \item \texttt{HPC\_QEMU/}
  \item \texttt{Zonal\_B/}
  \item \texttt{Telematics\_QEMU/}
  \item \texttt{PlantModel\_RestBus/}
\end{itemize}

\textbf{Support folders}
\begin{itemize}
  \item \texttt{log/}: captured run outputs.
  \item root README: run instructions and test case summary.
\end{itemize}

\textbf{Generated / external}
\begin{itemize}
  \item \texttt{build/} subfolders: generated by CMake.
  \item \texttt{install/}, \texttt{lib/}: SIL Kit SDK headers and binaries.
\end{itemize}

\section*{3. File-by-File and Function-by-Function Guide}

\subsection*{3.1 Zonal A}
\textbf{File: \texttt{Zonal\_A/ZonalA.cpp}}
\begin{itemize}
  \item \texttt{main}: creates participant \texttt{ZONAL1}, sets CAN receive callbacks and Ethernet send callbacks, starts lifecycle and worker thread.
\end{itemize}

\textbf{File: \texttt{Zonal\_A/include/Helper.hpp}}
\begin{itemize}
  \item \texttt{PrintHexFrame}: prints byte arrays in hex for debugging.
  \item \texttt{CAN\_FrameTransmitHandler}: logs CAN transmit status.
  \item \texttt{CAN\_FrameHandler}: receives CAN payload from PlantModel and stores it in global \texttt{Can\_payload}; sets send flag.
  \item \texttt{ETH\_FrameTransmitHandler}: Ethernet transmit status handling.
  \item \texttt{ETH\_FrameHandler}: Ethernet receive callback placeholder.
  \item \texttt{CreateFrame}: builds Ethernet raw frame with fixed MAC/header bytes plus received CAN payload.
  \item \texttt{ETH\_SendFrame}: sends the created Ethernet frame when new CAN data arrives.
\end{itemize}

\subsection*{3.2 HPC}
\textbf{File: \texttt{HPC\_QEMU/HPC.cpp}}
\begin{itemize}
  \item \texttt{Extract\_bits}: decomposes one input byte into 8 bits used for decision logic.
  \item \texttt{PrintHexFrame}, \texttt{operator<<}, \texttt{GetPayloadHexFromFrame}: utility helpers.
  \item \texttt{FrameHandler\_ZONAL\_A}: receives Ethernet from Zonal A, stores payload, extracts bits, sets hazard flag.
  \item \texttt{FrameHandle\_TELEMATIC}: receives trunk request data from Telematics and sets telematics flag.
  \item \texttt{CreateFrame\_1}: creates hazard frame for Zonal B.
  \item \texttt{CreateFrame\_2}: creates trunk frame for Zonal B.
  \item \texttt{CreateFrame\_2\_tele}: creates trunk response frame for Telematics.
  \item \texttt{Send\_HAZARD\_2\_ZONAL}: sends hazard state frames (including blink loop under active conditions).
  \item \texttt{Send\_TELEMATIC\_2\_ZONAL}: main trunk permission logic:
  \begin{itemize}
    \item open request + vehicle still = allow open,
    \item open request + moving = force closed,
    \item close request = close.
  \end{itemize}
  \item \texttt{Send\_HPC\_2\_TELEMATIC}: sends final trunk status back to Telematics.
  \item \texttt{main}: creates Eth1/Eth2/Eth3/Eth4 controllers, registers callbacks, runs two worker threads.
\end{itemize}

\subsection*{3.3 Zonal B}
\textbf{File: \texttt{Zonal\_B/src/ZonalB.cpp}}
\begin{itemize}
  \item \texttt{GetPayloadHexFromFrame}, \texttt{PrintHexFrame}: utility helpers.
  \item \texttt{ETH\_FrameHandler\_HPC\_HAZARD}: receives hazard Ethernet payload from HPC.
  \item \texttt{ETH\_FrameHandler\_HPC\_TRUNK}: receives trunk Ethernet payload from HPC.
  \item \texttt{SendFrame\_ZONALB\_2\_REST\_HAZARD}: converts hazard Ethernet payload to CAN ID \texttt{0x101} and sends on CAN2.
  \item \texttt{SendFrame\_ZONALB\_2\_REST\_TRUNK}: converts trunk Ethernet payload to CAN ID \texttt{0x103} and sends on CAN3.
  \item \texttt{FrameTransmitHandler\_TrunkSend}: logs trunk CAN tx status.
  \item \texttt{main}: participant setup, callbacks, periodic forwarding to PlantModel over CAN2/CAN3.
\end{itemize}

\subsection*{3.4 Telematics}
\textbf{File: \texttt{Telematics\_QEMU/Telematic\_ECU.cpp}}
\begin{itemize}
  \item \texttt{GetPayloadHexFromFrame}: extracts payload from Ethernet frame.
  \item \texttt{FrameHandler\_HPC}: receives HPC response and publishes user-facing MQTT status message.
  \item \texttt{main}: starts SIL Kit participant and two threads:
  \begin{itemize}
    \item MQTT listener thread for trunk commands,
    \item SIL Kit sender thread to forward trunk command frames to HPC.
  \end{itemize}
\end{itemize}

\textbf{File: \texttt{Telematics\_QEMU/include/Helper\_SIL\_Kit.hpp}}
\begin{itemize}
  \item \texttt{SendFrame}: sends \texttt{TRUNK\_FRAME} over Eth1 when MQTT flag is set.
  \item \texttt{FrameTransmitHandler}: Ethernet tx status logging.
\end{itemize}

\textbf{File: \texttt{Telematics\_QEMU/include/Telematics.hpp}}
\begin{itemize}
  \item \texttt{PrintHexFrame}: utility.
  \item \texttt{send\_udp\_message}: maps MQTT text command to binary \texttt{TRUNK\_FRAME} payload.
  \item \texttt{listen\_to\_mqtt}: subscribes using \texttt{mosquitto\_sub}, updates trunk frame and flag.
  \item \texttt{publish\_to\_mqtt}: publishes response using \texttt{mosquitto\_pub}.
\end{itemize}

\textbf{File: \texttt{Telematics\_QEMU/include/TrunkData.hpp}}
\begin{itemize}
  \item shared globals \texttt{TRUNK\_FRAME} and \texttt{CHECK\_MQTT}.
\end{itemize}

\subsection*{3.5 Plant Model RestBus}
\textbf{File: \texttt{PlantModel\_RestBus/RestBus.cpp}}
\begin{itemize}
  \item \texttt{main}: participant \texttt{PlantModel\_Restbus}, reads test case ID, sends initial stimulus, receives CAN feedback, coordinates test execution.
\end{itemize}

\textbf{File: \texttt{PlantModel\_RestBus/include/Helper.hpp}}
\begin{itemize}
  \item expected payload vectors: golden references for open/close/hazard patterns.
  \item \texttt{FrameHandler\_ReceiveHazard}: receives hazard CAN payload.
  \item \texttt{FrameTransmitHandler\_HazardSend}: logs outgoing hazard stimulus tx ack.
  \item \texttt{FrameHandler\_ReceiveTrunk}: compares received trunk payload with expected vectors, prints PASS/FAIL.
  \item \texttt{SendFrame\_ZonalA}: builds and sends CAN test stimulus (ID \texttt{0x401}) based on selected test case.
\end{itemize}

\subsection*{3.6 Build, Config, Scripts}
\begin{itemize}
  \item Each module has a \texttt{CMakeLists.txt}: compiles one executable, links SIL Kit and Threads.
  \item Each module has \texttt{config/participantConfig.yaml}: sets registry URI and logging.
  \item Each module has scripts:
  \begin{itemize}
    \item \texttt{first\_build.sh}: CMake configure + build,
    \item \texttt{start*.sh}: launches participant with expected args.
  \end{itemize}
\end{itemize}

\section*{4. Key Runtime Interfaces}
\textbf{CAN channels}
\begin{itemize}
  \item CAN1: PlantModel -> Zonal A stimulus.
  \item CAN2: Zonal B -> PlantModel hazard output.
  \item CAN3: Zonal B -> PlantModel trunk output.
\end{itemize}

\textbf{Ethernet channels}
\begin{itemize}
  \item Eth4: Zonal A <-> HPC.
  \item Eth1: Telematics <-> HPC.
  \item Eth2: HPC -> Zonal B hazard path.
  \item Eth3: HPC -> Zonal B trunk path.
\end{itemize}

\section*{5. MATLAB Replication Blueprint}
Use this order so implementation stays manageable:
\begin{enumerate}
  \item Build a pure MATLAB simulation (no sockets first).
  \item Represent each participant as a function/class with persistent state.
  \item Represent each bus as message queues (struct arrays).
  \item Keep payloads as \texttt{uint8} vectors.
  \item Implement helper functions:
  \begin{itemize}
    \item \texttt{createEthFrame(payload, srcMac, dstMac, prePayload)}
    \item \texttt{extractPayload(frame)}
    \item \texttt{zonalA\_step()}, \texttt{hpc\_step()}, \texttt{zonalB\_step()}, \texttt{telematics\_step()}, \texttt{plant\_step()}
  \end{itemize}
  \item Add test harness for TC1/TC2/TC3 and assert expected trunk results.
  \item After logic is stable, replace queue stubs with real transport (MQTT/UDP if needed).
\end{enumerate}

\section*{6. Important Engineering Notes}
When replicating in MATLAB, avoid carrying over fragile patterns from current C++ implementation:
\begin{itemize}
  \item Heavy use of globals and magic byte offsets (for example index 36/44) increases risk.
  \item Some synchronization lines use temporary mutex objects and do not protect shared data as intended.
  \item Some condition checks are likely unintended (for example string compare usage and always-true range check).
\end{itemize}

A cleaner MATLAB design should prefer:
\begin{itemize}
  \item explicit state structs,
  \item named signal fields over hardcoded offsets,
  \item deterministic state machine logic,
  \item test assertions for each decision path.
\end{itemize}

\section*{7. Main Files Reference}
\begin{itemize}
  \item \url{/home/kj/CES_2025_artifacts/Domain_Body_TestCase/ReadMe_Bodydomain_feature_testbench.md}
  \item \url{/home/kj/CES_2025_artifacts/Domain_Body_TestCase/Zonal_A/ZonalA.cpp}
  \item \url{/home/kj/CES_2025_artifacts/Domain_Body_TestCase/Zonal_A/include/Helper.hpp}
  \item \url{/home/kj/CES_2025_artifacts/Domain_Body_TestCase/HPC_QEMU/HPC.cpp}
  \item \url{/home/kj/CES_2025_artifacts/Domain_Body_TestCase/Zonal_B/src/ZonalB.cpp}
  \item \url{/home/kj/CES_2025_artifacts/Domain_Body_TestCase/Telematics_QEMU/Telematic_ECU.cpp}
  \item \url{/home/kj/CES_2025_artifacts/Domain_Body_TestCase/Telematics_QEMU/include/Helper_SIL_Kit.hpp}
  \item \url{/home/kj/CES_2025_artifacts/Domain_Body_TestCase/Telematics_QEMU/include/Telematics.hpp}
  \item \url{/home/kj/CES_2025_artifacts/Domain_Body_TestCase/Telematics_QEMU/include/TrunkData.hpp}
  \item \url{/home/kj/CES_2025_artifacts/Domain_Body_TestCase/PlantModel_RestBus/RestBus.cpp}
  \item \url{/home/kj/CES_2025_artifacts/Domain_Body_TestCase/PlantModel_RestBus/include/Helper.hpp}
\end{itemize}

\end{document}
